% !TeX root = ../../Book2.tex
% 纲要标题：### 5.4 范畴等价
% 纲要位置：计划纲要.md，第 970 行。

\section{范畴等价}
\label{b2:sec:0504}

上一节用积等构造的映射性质来辨认对象；若要比较两个范畴，则还须比较每一对对象之间的态射及其复合。选择坐标把线性映射写成矩阵，往往没有改变我们要研究的结构，却改变了对象的具体外形。范畴等价使这种关系成为可证明的陈述：它要求对象可以在同构意义下互相对应，也要求全部态射及其复合得到保留。

% 纲要内容（仅作写作计划，尚非教材正文）：
%
% * 全、忠实与本质满
% * 范畴等价、同构与泛性质
% * 坐标选择所得拟逆之间的自然同构
% * 为第16章 Morita 理论准备最低限度语言
%
% ---
%

\subsection{全、忠实与本质满}
\label{b2:subsec:0504:01}

\begin{definition}\label{b2:def:fully-faithful-essentially-surjective}
设 \(\mathcal C,\mathcal D\) 是范畴，\(F:\mathcal C\rightarrow\mathcal D\) 是函子。对 \(\mathcal C\) 中每对对象 \(X,Y\)，\(F\) 给出函数
\[
F_{X,Y}:\operatorname{Hom}_{\mathcal C}(X,Y)\longrightarrow
\operatorname{Hom}_{\mathcal D}\bigl(F(X),F(Y)\bigr),\qquad f\longmapsto F(f).
\]
若这些函数全部满射，称 \(F\) 为\term[quan-hanzi]{全函子}[full functor]；若全部单射，称为\term[zhongshi-hanzi]{忠实函子}[faithful functor]；若全部双射，称为\term[quanzhongshi-hanzi]{全忠实函子}[fully faithful functor]。若 \(\mathcal D\) 的每个对象都同构于某个 \(F(X)\)，称 \(F\) \term[benzhi-man]{本质满}[essentially surjective]。
\end{definition}

这三个条件检查的位置不同：
\begin{table}[htbp]
\centering
\caption{全、忠实与本质满的核对位置}\label{b2:tab:functor-equivalence-conditions}
\begin{tabularx}{\linewidth}{@{}p{0.19\linewidth}X X@{}}
\toprule
条件 & 对每个 Hom 映射的要求 & 对目标范畴对象的要求 \\
\midrule
忠实 & 单射 & 无 \\
全 & 满射 & 无 \\
本质满 & 无 & 每个对象同构于某个像对象 \\
\bottomrule
\end{tabularx}
\end{table}

“忠实”不合并固定起点、终点之间的不同态射，但不要求不同对象有不同的像；“全”要求目标中像对象之间的态射没有遗漏；“本质满”则要求目标对象在同构意义下没有遗漏。模到集合的忘却函子忠实：两个模同态若为同一个函数，就确实相等。当 \(R\ne0\) 时它不全，因为正则模底集合上的 \(r\mapsto r+1_R\) 不保持零元；零环的模均为单元素模，此时忘却函子反而全忠实。交换群到群的包含函子全忠实，因为交换群之间的群同态正是交换群同态；它不本质满，因为非交换群不能同构于交换群。

\ProseParagraphBreak
对偏序集之间的函子 \(F:P\rightarrow Q\)，函子性恰好是保序性。它总忠实，因为起点、终点固定后至多一个态射；它全当且仅当 \(F(x)\le F(y)\) 蕴含 \(x\le y\)。所以“保序”不等于“全”：把两元素反链送到两元素链的对应位置，是忠实而不全的函子。

\begin{proposition}\label{b2:prop:fully-faithful-reflects-isomorphism}
设 \(F:\mathcal C\to\mathcal D\) 是范畴间的函子。\(F\) 把 \(\mathcal C\) 中的同构送到 \(\mathcal D\) 中的同构。若 \(F\) 全忠实，则对 \(\mathcal C\) 中的态射 \(f:X\to Y\)，\(F(f)\) 可逆蕴含 \(f\) 可逆；对对象 \(X,Y\)，\(F(X)\cong F(Y)\) 蕴含 \(X\cong Y\)。
\end{proposition}
\begin{proof}
若 \(g=f^{-1}\)，函子公理给出 \(F(g)F(f)=F(gf)=\operatorname{id}\) 及反序的同样等式，故 \(F(g)\) 为逆。

现设 \(F\) 全忠实。若 \(F(f)\) 的逆为 \(v:F(Y)\rightarrow F(X)\)，要证明 \(f\) 可逆，须先把 \(v\) 提升为候选逆，再把两条逆关系从像中反映回来。全性给出 \(g:Y\rightarrow X\) 使 \(F(g)=v\)，于是
\[
\begin{aligned}
F(gf)&=vF(f)=\operatorname{id}_{F(X)}=F(\operatorname{id}_X),\\
F(fg)&=F(f)v=\operatorname{id}_{F(Y)}=F(\operatorname{id}_Y).
\end{aligned}
\]
忠实性分别给出 \(gf=\operatorname{id}_X\)、\(fg=\operatorname{id}_Y\)，故 \(g=f^{-1}\)。若只给定 \(a:F(X)\rightarrow F(Y)\) 为同构，先用全性把 \(a\) 提升为 \(f\)，再用刚证明的结论即可。
\end{proof}

\subsection{范畴等价、同构与泛性质}
\label{b2:subsec:0504:02}

\begin{definition}\label{b2:def:category-equivalence}
设 \(\mathcal C,\mathcal D\) 是范畴，\(F:\mathcal C\rightarrow\mathcal D\) 是函子。若存在函子 \(G:\mathcal D\rightarrow\mathcal C\) 及自然同构
\[
\eta:\operatorname{Id}_{\mathcal C}\Rightarrow GF,\qquad
\varepsilon:FG\Rightarrow\operatorname{Id}_{\mathcal D},
\]
则称 \(F\) 为\term[fanchou-dengjia]{范畴等价}[equivalence of categories]，\(G\) 为其\term[zhunni]{拟逆}[quasi-inverse]，写作 \(\mathcal C\simeq\mathcal D\)。若可要求 \(GF=\operatorname{Id}_{\mathcal C}\)、\(FG=\operatorname{Id}_{\mathcal D}\) 作为函子严格相等，则称 \(F\) 为\term[fanchou-tonggou]{范畴同构}[isomorphism of categories]。
\end{definition}
\symidx{\(\mathcal C\simeq\mathcal D\)：范畴等价}

定义中的 \(\eta,\varepsilon\) 都是自然同构，因而也可以取它们的逆，用相反方向表达同一个比较。重要的是 \(GF\) 与 \(\operatorname{Id}_{\mathcal C}\)、\(FG\) 与 \(\operatorname{Id}_{\mathcal D}\) 分别自然同构。范畴同构在对象上给出真正的双射；等价则允许多个互相同构的对象合并到一个代表。比如范畴 \(\mathcal E\) 有两个对象 \(a,b\)，每对对象之间恰有一个态射，复合由唯一性确定。把它们都送到仅有一个对象及恒等态射的范畴 \(\mathbf 1\)，并反向选取 \(a\)，便给出等价：每个对象到 \(a\) 的唯一态射组成所需自然同构。但这两个范畴对象数不同，故不可能同构。

\ProseParagraphBreak
构造拟逆时，本质满使目标对象能找到同构的像对象；选定这些代表后，就要把目标中的态射经同构搬到两个像对象之间。全性保证搬来的态射有原像，忠实性保证原像唯一，并将恒等和复合的等式反映回原范畴。下面把这些要求分别落实到对象和态射的构造中。

\begin{theorem}[范畴等价的判据]\label{b2:thm:equivalence-criterion}
设 \(\mathcal C,\mathcal D\) 是范畴，\(F:\mathcal C\rightarrow\mathcal D\) 是函子。
\begin{enumerate}
\item 若 \(F\) 是范畴等价，则它全忠实且本质满。
\item 若 \(F\) 全忠实，并已为每个 \(Y\in\mathcal D\) 同时指定对象 \(G(Y)\in\mathcal C\) 及同构 \(\varepsilon_Y:F(G(Y))\to Y\)，则 \(F\) 是范畴等价。
\end{enumerate}
\end{theorem}
\begin{proof}
先设 \(G,\eta,\varepsilon\) 给出等价。对每个 \(Y\in\mathcal D\)，同构 \(\varepsilon_Y:F\bigl(G(Y)\bigr)\rightarrow Y\) 证明本质满。若 \(f,g:X\rightarrow X'\) 满足 \(F(f)=F(g)\)，\(\eta\) 的自然性给出
\[
\eta_{X'}f=GF(f)\eta_X=GF(g)\eta_X=\eta_{X'}g.
\]
消去同构 \(\eta_{X'}\)，得 \(f=g\)，所以 \(F\) 忠实。为了随后验证提升来的态射确实具有所需的 \(F\) 像，还要用到 \(G\) 的忠实性。若 \(u,v:Y\rightarrow Y'\) 且 \(G(u)=G(v)\)，\(\varepsilon\) 的自然性给出
\[
u\varepsilon_Y=\varepsilon_{Y'}FG(u)
=\varepsilon_{Y'}FG(v)=v\varepsilon_Y.
\]
消去 \(\varepsilon_Y\) 得 \(u=v\)，所以 \(G\) 也忠实。

给定 \(a:F(X)\rightarrow F(X')\)，先用 \(G\) 得到 \(G(a):GF(X)\rightarrow GF(X')\)，再经 \(\eta_X\) 与 \(\eta_{X'}^{-1}\) 搬回 \(X,X'\) 之间，得到候选态射
\[
f=\eta_{X'}^{-1}G(a)\eta_X:X\longrightarrow X'.
\]
由 \(\eta\) 对 \(f\) 的自然性，\(GF(f)\eta_X=\eta_{X'}f=G(a)\eta_X\)。消去 \(\eta_X\)，得 \(G\bigl(F(f)\bigr)=G(a)\)；再由 \(G\) 忠实得 \(F(f)=a\)。故 \(F\) 全。

反过来，设 \(F\) 全忠实，并已同时指定所述 \(G(Y)\) 与 \(\varepsilon_Y\)。对 \(a:Y\rightarrow Y'\)，要定义 \(G(a)\)，应先要求它在 \(F\) 下的像与这些代表同构相容，即下图交换：
\[
\begin{tikzcd}[column sep=3em,row sep=2em]
F(G(Y))\arrow[r,dashed,"F(G(a))"]\arrow[d,"\varepsilon_Y"']&
F(G(Y'))\arrow[d,"\varepsilon_{Y'}"]\\
Y\arrow[r,"a"']&Y'.
\end{tikzcd}
\]
因而中间态射必须先经 \(\varepsilon_Y\) 到 \(Y\)，再施加 \(a\)，最后经 \(\varepsilon_{Y'}^{-1}\) 回到 \(F(G(Y'))\)。全性使这条搬来的态射有原像，忠实性使原像唯一，所以存在唯一 \(G(a):G(Y)\rightarrow G(Y')\) 满足
\begin{equation}\label{b2:eq:quasi-inverse-on-arrows}
F\bigl(G(a)\bigr)=\varepsilon_{Y'}^{-1}a\varepsilon_Y.
\end{equation}
当 \(a=\operatorname{id}_Y\) 时，右边为恒等态射，忠实性给出 \(G(\operatorname{id}_Y)=\operatorname{id}_{G(Y)}\)。对 \(a:Y\rightarrow Y'\)、\(b:Y'\rightarrow Y''\)，有
\[
\begin{aligned}
F\bigl(G(b)G(a)\bigr)
&=\varepsilon_{Y''}^{-1}b\varepsilon_{Y'}\varepsilon_{Y'}^{-1}a\varepsilon_Y\\
&=\varepsilon_{Y''}^{-1}ba\varepsilon_Y=F\bigl(G(ba)\bigr).
\end{aligned}
\]
忠实性给出 \(G(b)G(a)=G(ba)\)，故 \(G\) 是函子。式\eqref{b2:eq:quasi-inverse-on-arrows}恰好说 \(\varepsilon\) 是自然同构。

还需把 \(X\) 与 \(GF(X)\) 比较。将 \(Y=F(X)\) 代入已选的代表同构，得到 \(\varepsilon_{F(X)}:F(GF(X))\rightarrow F(X)\)。它的逆方向恰好是 \(F(X)\rightarrow F(GF(X))\)，因此用全忠实性唯一地提升这个逆，得到 \(\eta_X:X\rightarrow GF(X)\)，满足
\[
F(\eta_X)=\varepsilon_{F(X)}^{-1}.
\]
右端可逆，\cref{b2:prop:fully-faithful-reflects-isomorphism}说明 \(\eta_X\) 可逆。对 \(f:X\rightarrow X'\)，式\eqref{b2:eq:quasi-inverse-on-arrows}用于 \(F(f)\) 后给出
\[
\begin{aligned}
F\bigl(GF(f)\eta_X\bigr)
&=\varepsilon_{F(X')}^{-1}F(f)\varepsilon_{F(X)}\varepsilon_{F(X)}^{-1}\\
&=\varepsilon_{F(X')}^{-1}F(f)
=F\bigl(\eta_{X'}f\bigr).
\end{aligned}
\]
由忠实性，\(GF(f)\eta_X=\eta_{X'}f\)，所以 \(\eta\) 自然，所需等价得证。

\end{proof}

若 \(\mathcal C,\mathcal D\) 都是小范畴，本质满给出的候选对 \((X,a)\)、\(a:F(X)\xrightarrow{\sim}Y\) 组成非空集合，并由集合 \(\operatorname{Ob}(\mathcal D)\) 编号；选择公理因此能同时选出定理第二项所需的代表。故对小范畴，等价恰好等于全忠实且本质满。对任意大范畴，逐个对象的存在性还不自动给出这种统一选择，因此定理将已指定的同构代表作为数据。若按\secref{b2:sec:0501}的大小约定，两个范畴在某个更大宇宙中成为小范畴，则可在该宇宙中作出选择。这项判据及拟逆构造可对照 Riehl\cite[Theorem 1.5.9, pp. 33--34]{Riehl2016CategoryTheory}。

\ProseParagraphBreak
设 \(k\) 是域，定义矩阵范畴 \(\mathbf{Mat}_k\)：对象为非负整数，\(n\) 到 \(m\) 的态射是 \(m\times n\) 矩阵，复合为矩阵乘法，恒等态射为单位矩阵；零尺寸按唯一的空矩阵处理。函子
\[
J:\mathbf{Mat}_k\longrightarrow\mathbf{Vect}_k^{\mathrm{fd}},\qquad
n\longmapsto k^n,\quad A\longmapsto(v\mapsto Av)
\]
\symidx{\(\mathbf{Mat}_k,\mathbf{Vect}_k^{\mathrm{fd}}\)：矩阵范畴与有限维向量空间范畴}
全忠实，因为标准基下每个线性映射恰由一个矩阵表示；它本质满，因为有限维空间有基，因而同构于 \(k^n\)。这里固定一个包含 \(k\) 的宇宙，只考虑其中的有限维向量空间；在更大的宇宙中，两范畴都小，因而可用上述判据得到等价。

\ProseParagraphBreak
具体选定每个有限维 \(V\) 的坐标同构 \(c_V:k^{\dim_k V}\rightarrow V\)，定义拟逆
\[
H:\mathbf{Vect}_k^{\mathrm{fd}}\longrightarrow\mathbf{Mat}_k,
\qquad V\longmapsto\dim_kV,
\]
并把 \(f:V\rightarrow W\) 送到 \(c_W^{-1}fc_V\) 的矩阵。复合时中间的坐标同构相消，给出 \(H(gf)=H(g)H(f)\)，且 \(H(\operatorname{id}_V)=I_{\dim_kV}\)，所以 \(H\) 是函子。这里 \(c_V\) 的源为 \(JH(V)\)，目标为 \(V\)，所以它们组成的比较是 \(c:JH\Rightarrow\operatorname{Id}_{\mathbf{Vect}_k^{\mathrm{fd}}}\)。等式 \(fc_V=c_WJH(f)\) 给出自然性方块：
\[
\begin{tikzcd}[column sep=3em,row sep=2em]
JH(V)\arrow[r,"c_V"]\arrow[d,"JH(f)"']&V\arrow[d,"f"]\\
JH(W)\arrow[r,"c_W"']&W.
\end{tikzcd}
\]
每个 \(c_V\) 都可逆，故\cref{b2:prop:natural-isomorphism-components}使 \(c\) 成为自然同构。在矩阵一侧，\(HJ(n)=n\)，而标准基下 \(c_{k^n}^{-1}\) 的矩阵 \(\eta_n:n\to HJ(n)\) 满足
\[
\begin{aligned}
HJ(A)\eta_n&=c_{k^m}^{-1}Ac_{k^n}c_{k^n}^{-1}\\
&=c_{k^m}^{-1}A=\eta_mA
\qquad(A:n\to m).
\end{aligned}
\]
因此 \((\eta_n)\) 给出自然同构 \(\operatorname{Id}_{\mathbf{Mat}_k}\Rightarrow HJ\)。拟逆 \(H\) 及这些比较依赖于坐标选择，\(J\) 则只使用标准基。维数只描述对象；\(H\) 还必须记录同态的矩阵，才能保持态射与复合。

\begin{proposition}\label{b2:prop:coordinate-choice-natural}
设 \(k\) 是域，在有限维 \(k\) 向量空间范畴 \(\mathbf{Vect}_k^{\mathrm{fd}}\) 中，为每个对象 \(V\) 指定两族坐标同构 \(c_V,d_V:k^{\dim_kV}\to V\)。令 \(H_c,H_d:\mathbf{Vect}_k^{\mathrm{fd}}\to\mathbf{Mat}_k\) 在对象上均取 \(V\mapsto\dim_kV\)，在态射 \(f:V\to W\) 上分别取 \(c_W^{-1}fc_V\)、\(d_W^{-1}fd_V\) 的矩阵。则标准基下 \(d_V^{-1}c_V\) 的矩阵 \(P_V\) 组成自然同构 \(P:H_c\Rightarrow H_d\)。
\end{proposition}
\begin{proof}
对每个 \(f:V\to W\)，矩阵乘法给出
\[
H_d(f)P_V=d_W^{-1}fd_Vd_V^{-1}c_V
=d_W^{-1}fc_V=P_WH_c(f).
\]
这就是 \(P\) 的自然性。每个 \(P_V\) 都可逆，其逆为 \(c_V^{-1}d_V\) 的矩阵；由\cref{b2:prop:natural-isomorphism-components}，\(P\) 是自然同构。因而不同坐标选择虽然可能产生不同的拟逆，它们由这些换基矩阵自然地联系起来。
\end{proof}

泛性质只规定相容态射的存在与唯一性。等价的全忠实性保留每个 Hom 集合及复合，本质满使任意目标测试对象都能换成同构的像对象，因此这些存在与唯一性要求可以在两个范畴之间传递。

\begin{proposition}\label{b2:prop:equivalence-preserves-universal-objects}
设 \(F:\mathcal C\to\mathcal D\) 是范畴等价。\(F\) 把 \(\mathcal C\) 的初始对象、终对象分别送到 \(\mathcal D\) 的初始对象、终对象。对任意由集合 \(I\) 索引的对象族 \((X_i)_{i\in I}\)，若其积或余积存在，\(F\) 还把该积或余积连同指定映射送到 \((F(X_i))_{i\in I}\) 的积或余积。
\end{proposition}
\begin{proof}
由\cref{b2:thm:equivalence-criterion}的正向结论，\(F\) 全忠实且本质满。若 \(I_0\) 是初始对象，对任意 \(Y\in\mathcal D\)，选同构 \(a:F(T)\rightarrow Y\)。全忠实性与后复合 \(a\) 给出双射
\[
\operatorname{Hom}_{\mathcal C}(I_0,T)
\longrightarrow\operatorname{Hom}_{\mathcal D}\bigl(F(I_0),Y\bigr).
\]
左边恰一个元素，故 \(F(I_0)\) 初始。若 \(Z\) 是终对象，预复合 \(a^{-1}\) 把 \(\operatorname{Hom}_{\mathcal C}(T,Z)\) 双射到 \(\operatorname{Hom}_{\mathcal D}\bigl(Y,F(Z)\bigr)\)，故 \(F(Z)\) 是终对象。

设 \((P,p_i)\) 为 \((X_i)_{i\in I}\) 的积。给定 \(b_i:Y\rightarrow F(X_i)\)，仍选 \(a:F(T)\rightarrow Y\)。全忠实性给出唯一 \(c_i:T\rightarrow X_i\) 使 \(F(c_i)=b_i a\)。积的泛性质给出唯一 \(c:T\rightarrow P\)，使 \(p_i c=c_i\)。于是 \(b=F(c)a^{-1}:Y\rightarrow F(P)\) 满足 \(F(p_i)b=b_i\)。若另一态射 \(b'\) 也满足，先把 \(b'a:F(T)\rightarrow F(P)\) 唯一提升为 \(c':T\rightarrow P\)；忠实性给出 \(p_ic'=c_i\)，积的唯一性给出 \(c'=c\)，所以 \(b'=b\)。

对余积 \((Q,j_i)\)，给定 \(b_i:F(X_i)\rightarrow Y\)，把 \(a^{-1}b_i\) 唯一提升为 \(c_i:X_i\rightarrow T\)，由余积得到 \(c:Q\rightarrow T\)，再令 \(b=aF(c)\)。若 \(b'F(j_i)=b_i\)，把 \(a^{-1}b'\) 提升后，用余积的唯一性得到同一个 \(c\)。故余积也被保留。
\end{proof}

这里保留积的结论，是 \(F(P)\) 连同 \(F(p_i)\) 满足目标范畴的积泛性质，因而与任何另一个积存在唯一相容同构；不要求 \(F(P)\) 与某个底集合的坐标直积严格相等。余积也必须连同指定嵌入一起比较，只有对象的抽象同构类型不足以检验这种相容性。

\subsection{Morita 理论的范畴语言}
\label{b2:subsec:0504:03}

若两个环 \(R,S\) 同构，它们的模论只相差标量的名称。设 \(\varphi:R\rightarrow S\) 为环同构，对左 \(R\) 模 \(M\) 定义左 \(S\) 作用
\[
s\cdot m=\varphi^{-1}(s)m.
\]
在模同态上保持原映射，得到 \(R\text{-}\mathbf{Mod}\rightarrow S\text{-}\mathbf{Mod}\) 的函子。以 \(\varphi\) 替代 \(\varphi^{-1}\) 作反向构造，两次作用恢复原作用，两个函子严格互逆。

\ProseParagraphBreak
如果只要求模范畴等价，是否仍能恢复环同构？\chapref{b2:ch:16}将研究这一问题，并给出比环同构更宽的 \term[Morita-dengjia]{Morita 等价}[Morita equivalence]：它指 \(R\text{-}\mathbf{Mod}\) 与 \(S\text{-}\mathbf{Mod}\) 的范畴等价。那里还将证明矩阵环与原环的具体等价。

\ProseParagraphBreak
由\cref{b2:prop:equivalence-preserves-universal-objects}，模范畴等价保留零模、直和及直积所表达的映射性质，但正则模不是这个等价自动识别的指定对象。若等价把 \({}_RR\) 送到某个 \(S\) 模 \(P\)，它所对应的是 \(P\) 的自同态及其复合，并不保证 \(P\cong{}_SS\)。\cref{b2:prop:regular-hom-endomorphism}给出 \(\operatorname{End}_R({}_RR)\cong R^{\mathrm{op}}\)，其中反环来自右乘与复合次序；要在另一边用同一公式识别 \(S^{\mathrm{op}}\)，还须知道所取对象确实是正则模。因此，模范畴等价本身不允许直接把两个正则模的自同态环等同起来。\chapref{b2:ch:16}将研究哪些模能够代替正则模来描述整个模范畴。
